\documentclass[a4paper,11pt]{article}
%% ── Paquets externes (noyau pour atome) ───────────────
\usepackage[utf8]{inputenc}
\usepackage[T1]{fontenc}
\usepackage{lmodern}
\usepackage[french]{babel}
\usepackage{geometry}
\usepackage{amsmath}
\usepackage{amssymb}
\usepackage{etoolbox}
\usepackage{xkeyval}
\usepackage{xifthen}
\usepackage{environ}
\usepackage{forloop}
\usepackage{xcolor}
\usepackage{ninecolors}
\usepackage{multirow}
\usepackage{colortbl}
\usepackage{tabularray}
\usepackage{multicol}
\usepackage[inline]{enumitem}
\usepackage{graphicx}
\usepackage{adjustbox}
\usepackage{tikz}
\usepackage[theorems,breakable,skins]{tcolorbox}
\usepackage{pifont}

%% ── Marges (alignées sur les livrets) ───────────────
\newgeometry{left=35pt,right=35pt,top=65pt,bottom=70pt}

%% ── Paquets externes additionnels (selon l'atome) ─────
\usepackage{seqenseigne-style6e}

%% ── Définitions du paquet seqenseigne (inlinées) ─────
%%    Fermeture transitive des macros effectivement utilisées
%%    par l'atome — au lieu de \usepackage{seqenseigne}
%%    qui chargerait ~110 paquets.
\makeatletter
\makeatother

%% ── Macros transverses ───────────────────────────────
\providecommand{\acompleter}[1]{\;\underline{\phantom{#1}}\;}

\begin{document}

\seqInitCorrigesEval
\seqTitreEval[
    titre={Évaluation --- Nombres},
    etab={\ldots},
    classe={\dots}
]


%% Barème global
\begin{seqEvalBareme}
    Exercice I & 5 points \\
    Exercice II & 3 points \\
    Exercice III & 2 points \\
\end{seqEvalBareme}


%% Tableau des objectifs et niveaux de maîtrise
\begin{seqEvalObjectifs}
    S01.Obj. 02 & Utiliser la valeur des chiffres
        & \ding{113} & \ding{113} & \ding{113} & \ding{113} \\
    S01.Obj. 03 & Écrire un grand nombre
        & \ding{113} & \ding{113} & \ding{113} & \ding{113} \\
\end{seqEvalObjectifs}


\begin{seqEvalExercice}[titre={Valeur des chiffres}, bareme={5 points}]
Sujet A1 \seqACompleter{3}.
\seqEvalCorrigeExo{%
\emph{Corrigé non rédigé.}
}
\end{seqEvalExercice}


\begin{seqEvalExercice}[titre={Grands nombres}, bareme={3 points}]
Sujet A2.
\seqEvalCorrigeExo{%
Explication A2.
}
\end{seqEvalExercice}


\begin{seqEvalExercice}[titre={Bonus : Qui suis-je ?}, bareme={2 points}]
Bonus A.
\seqEvalCorrigeExo{%
\emph{Corrigé non rédigé.}
}
\end{seqEvalExercice}


\seqEvalAfficheCorriges
\end{document}
